\documentclass{article}
\usepackage{amsmath, amssymb, amsthm}

\setlength{\parindent}{0pt}

\newtheorem{claim}{Claim}

\begin{document}

Suppose $n$ balls are thrown randomly into $n$ boxes, so each ball lands in each box with uniform probability. 
Also, suppose the outcome of each throw is independent of all the other throws.

\begin{claim}
    Let $X_i$ be an indicator random variable whose value is $1$ if box $i$ is empty and $0$ otherwise. The 
    closed form expression for the probability distribution of $X_i$ is 
    \[
    \Pr\{X_i = x\} =
    \begin{cases}
    \left(1 - \frac{1}{n}\right)^n, & x = 1, \\[6pt]
    1 - \left(1 - \frac{1}{n}\right)^n, & x = 0.
    \end{cases}
    \]
    and $X_1, X_2, \dots, X_n$ are dependent random variables.
\end{claim}

\begin{proof}
    Since each ball lands in each box with uniform probability, the probability that a ball lands in box $i$ is
    \[
    \Pr\{\text{ball lands in box } i\} = \frac{1}{n}
    \]
    Taking the complement,
    \[
    \Pr\{\text{ball does not land in box } i\} = 1 - \frac{1}{n}
    \]
    Since the throws are independent, the probability that box $i$ is empty is 
    \[
    \Pr\{\text{box } i \text{ is empty}\} = \left(1 - \frac{1}{n}\right)^n
    \]
    Hence
    \[
    \Pr\{X_i = 1\} = \left(1 - \frac{1}{n}\right)^n
    \]
    and 
    \[
    \Pr\{X_i = 0\} = 1 - \left(1 - \frac{1}{n}\right)^n
    \]

    \medskip

    It remains to verify whether the random variables $X_1, X_2, \dots, X_n$ are independent. If $X_1$ and $X_2$ 
    were independent, then
    \[
    \Pr\{X_1 = 1, X_2 = 1\} = \Pr\{X_1 = 1\}\Pr\{X_2 = 1\}
    \]
    We now determine whether this equality holds. Using the prior result, we have
    \[
    \Pr\{X_1 = 1\}\Pr\{X_2 = 1\} = \left(\left(1 - \frac{1}{n}\right)^n\right)^2
    \]
    To find $\Pr\{X_1 = 1, X_2 = 1\}$, we first calculate the probability that a ball 
    lands in either box $1$ or $2$.
    \[
    \Pr\{\text{ball lands in either box } 1 \text{ or } 2\} = \frac{1}{n} + \frac{1}{n} = \frac{2}{n}
    \]
    Taking the complement,
    \[
    \Pr\{\text{ball lands in neither box } 1 \text{ nor } 2\} = 1 - \frac{2}{n}
    \]
    Since the throws are independent, the probability that both boxes $1$ and $2$ are 
    empty is
    \[
    \Pr\{\text{box } 1 \text{ and } 2 \text{ are empty}\} = \left(1 - \frac{2}{n}\right)^n
    \]
    Hence
    \[
    \Pr\{X_1 = 1, X_2 = 1\} = \left(1 - \frac{2}{n}\right)^n
    \]
    However,
    \[
    \left(1 - \frac{2}{n}\right)^n \neq \left(\left(1 - \frac{1}{n}\right)^n\right)^2
    \]
    in general. For example, when $n = 2$,
    \[
    \left(1 - \frac{2}{2}\right)^2 = 0 \quad \text{while} \quad 
    \left(\left(1 - \frac{1}{2}\right)^2\right)^2 = \frac{1}{16}
    \]
    Hence the random variables $X_1, X_2, \dots, X_n$ are not independent.

    \medskip

    Hence the probability distribution of $X_i$ is
    \[
    \Pr\{X_i = x\} =
    \begin{cases}
    \left(1 - \frac{1}{n}\right)^n, & x = 1, \\[6pt]
    1 - \left(1 - \frac{1}{n}\right)^n, & x = 0.
    \end{cases}
    \]
    and $X_1, X_2, \dots, X_n$ are dependent random variables.
\end{proof}

\end{document}
